% !TEX root = ../metatime_monograph.tex
\chapter{The Gell-Mann--Okubo Mass Formula}
\label{ch:gmo_formula}

\section{Historical Background}

The Gell-Mann--Okubo (GMO) mass formula~\cite{gellmann1961,okubo1962}
was one of the early triumphs of SU(3) flavour symmetry.  It relates the
masses of baryons within the same SU(3) multiplet through a linear
expression in hypercharge \(Y\) and isospin \(I\):

\begin{equation}
M = a + bY + c\left[I(I+1) - \frac{Y^2}{4}\right],
\label{eq:gmo_standard}
\end{equation}

where \(a\), \(b\), and \(c\) are constants determined by the symmetry
breaking pattern.

\section{Derivation from SU(3) Breaking}

The SU(3) flavour symmetry is broken by the strange quark mass, which
transforms as the eighth component of an octet.  To first order in the
breaking, the mass operator is:

\[
M = M_0 + M_8 T_8,
\]

where \(T_8\) is the eighth Gell-Mann matrix.  In a given SU(3) multiplet,
the eigenvalues of \(T_8\) are linear combinations of \(Y\) and \(I(I+1)\),
leading to the GMO form.

For the baryon octet, this gives:

\[
M = M_0 + \alpha + \beta Y + \gamma\left[I(I+1) - \frac{Y^2}{4}\right],
\label{eq:gmo_metatime}
\]

where \(M_0\) is the common mass scale and \(\alpha, \beta, \gamma\) are
SU(3)-breaking coefficients.

\section{GMO Relations}

The GMO formula implies a mass relation within the octet.  Using the
assignments:

\[
\begin{array}{lccc}
\text{Particle} & Y & I & I(I+1) - Y^2/4 \\
\hline
N (p,n) & +1 & 1/2 & 1/2 \\
\Lambda & 0 & 0 & 0 \\
\Sigma & 0 & 1 & 2 \\
\Xi & -1 & 1/2 & 1/2
\end{array}
\]

we obtain:

\[
\begin{aligned}
M_N &= M_0 + \alpha + \beta + \gamma/2, \\
M_\Lambda &= M_0 + \alpha, \\
M_\Sigma &= M_0 + \alpha + 2\gamma, \\
M_\Xi &= M_0 + \alpha - \beta + \gamma/2.
\end{aligned}
\]

From these, the GMO relation follows:

\[
M_N + M_\Xi = \frac{3M_\Lambda + M_\Sigma}{2}.
\]

\section{GMO in the Metatime Framework}

In the Metatime framework, the coefficients \(\alpha, \beta, \gamma\) are not
free parameters but are derived from the L-constants and quark primes, as we
show in Chapter~\ref{ch:octet_coefficients}.  The common mass \(M_0\) is
derived from the exponential mass formula (Chapter~\ref{ch:octet_m0}).
Thus the entire baryon octet spectrum is determined by the same constants
that govern all other sectors.
